%% This document is created by 
%%  Dr. Putu Harry Gunawan
%% Edited by
%%  Wisnu Adi Nurcahyo

\documentclass[a4paper,12pt,oneside]{book}
\usepackage[utf8]{inputenc}
\usepackage{sectsty}
\usepackage{graphicx}
\usepackage{epstopdf}
\usepackage{algorithm}
\usepackage{algpseudocode}
\usepackage{array}
\usepackage[table]{xcolor}
\usepackage{anysize}
\usepackage{amsmath}
\usepackage{amssymb}
\usepackage[bahasa]{babel}
\usepackage{indentfirst} %Spasi untuk paragraf pertama
\usepackage{geometry}
\usepackage{multirow}% http://ctan.org/pkg/multirow
\usepackage{hhline}% http://ctan.org/pkg/hhline
\marginsize{4cm}{3cm}{3cm}{3cm} %{left}{right}{top}{bottom}
\usepackage[compact]{titlesec} 
\usepackage{etoolbox}
\usepackage{natbib} %for Harvard syle
\usepackage[bookmarks,hypertexnames=false,debug]{hyperref}
%\usepackage[pageanchor]{hyperref}
\usepackage{bookmark}
\usepackage{dirtytalk}
% \usepackage{ccg-latex}
% \usepackage{booktabs}

% CCG
\usepackage[TABcline]{tabstackengine}

\def\tbs{\,\textbackslash}
%

\makeatletter
\patchcmd{\ttlh@hang}{\parindent\z@}{\parindent\z@\leavevmode}{}{}
\patchcmd{\ttlh@hang}{\noindent}{}{}{}

% CCG
\newcommand{\so}[1]{\ensuremath{\mbox{\it #1}^\prime}}
\newcommand\CCGFA{%
  \rlap{\kern.5\dimexpr\tabbed@gap\relax\llap{\fboxsep=0pt\colorbox{white}{$>$}}}}
\newcommand\CCGBA{%
  \rlap{\kern.5\dimexpr\tabbed@gap\relax\llap{\fboxsep=0pt\colorbox{white}{$<$}}}}
\newcommand\CCGCOOR{%
  \rlap{\kern.5\dimexpr\tabbed@gap\relax\llap{\fboxsep=0pt\colorbox{white}{\footnotesize $<\&>$}}}}
\newcommand\CCGTR{%
  \rlap{\kern.5\dimexpr\tabbed@gap\relax\llap{\fboxsep=0pt\colorbox{white}{\footnotesize $>\text{T}$}}}}
%
\makeatother

% CCG
\setlength\TABruleshift{\dimexpr.5\ht\strutbox-.5\dp\strutbox}
\setstacktabbedgap{1ex}
\TABstackTextstyle{\itshape}
\setstackgap{S}{2pt}
%

\chapterfont{\centering}
\newcommand{\bigsize}{\fontsize{16pt}{14pt}\selectfont}
\chapterfont{\centering\bigsize\bfseries}
\sectionfont{\large\bfseries}
\usepackage{tikz}
\usetikzlibrary{shapes.geometric, arrows}
%\renewcommand{\chaptertitle}{BAB}
\renewcommand{\thechapter}{\Roman{chapter}}
\renewcommand\thesection{\arabic{chapter}.\arabic{section}}
\renewcommand\thesubsection{\thesection.\arabic{subsection}}
\renewcommand{\theequation}{\arabic{chapter}.\arabic{equation}}
\renewcommand{\thefigure}{\arabic{chapter}.\arabic{figure}}
\renewcommand{\thetable}{\arabic{chapter}.\arabic{table}}

\renewcommand\bibname{Daftar Pustaka}
\addto{\captionsbahasa}{\renewcommand{\bibname}{Daftar Pustaka}}
\usepackage{fancyhdr}
\pagestyle{fancy}
\lhead{}
\chead{}
\rhead{}
\lfoot{}
\cfoot{\thepage}
\rfoot{}
\renewcommand{\headrulewidth}{0pt}

\makeatletter

%%%%%%%%%%%%%%%%%%%%%%%%%%%%%%%%%%%%%%%%%%%%%%%%%%%%%%%%%%%%
%
%  Berikut adalah data-data yang wajib diisi oleh mahasiswa
%
%%%%%%%%%%%%%%%%%%%%%%%%%%%%%%%%%%%%%%%%%%%%%%%%%%%%%%%%%%%%

% Judul dalam bahasa Indonesia
\title{
  Membangun sebuah Combinatory Categorial Grammar (CCG) Supertagger
  Berbasis Maximum Entropy untuk Bahasa Indonesia
}\let\Title\@title

% Judul dalam bahasa Inggris
\newcommand{\EngTitle}{
  Building a Combinatory Categorial Grammar (CCG) Supertagger
  Based on the Maximum Entropy for Bahasa Indonesia
}

% Nama Penulis
\author{
  Wisnu Adi Nurcahyo
}\let\Author\@author

% NIM Penulis
\newcommand{\NIM}{1301160479}

% Prodi Penulis
\newcommand{\Prodi}{Informatika}

% Kelompok Keahlian TA
\newcommand{\KK}{Intelligence, Computing, and Multimedia (ICM)}

% Gelar
\newcommand{\Gelar}{Sains Komputasi}

% Tahun TA
\date{2019}\let\Date\@date

% Kelas TA
\newcommand{\Kelas}{NLP}

% Tanggal Pengesahan
\newcommand{\Tanggal}{13}

% Bulan Pengesahan
\newcommand{\Bulan}{November}

% Nama Pembimbing Satu
\newcommand{\PembimbingSatu}{Dr. Ade Romadhony, S.T., M.T.}

% NIP Pembimbing Satu
\newcommand{\NIPPembimbingSatu}{06840042}

% Nama Pembimbing Dua
\newcommand{\PembimbingDua}{Dr. Calon Pembimbing 2, M.Kom}

% NIP Pembimbing Dua
\newcommand{\NIPPembimbingDua}{123456}

% Nama Kaprodi
\newcommand{\Kaprodi}{Niken Dwi Wahyu Cahyani, S.T., M.Kom. PhD}

% NIP Kaprodi
\newcommand{\NIPKaprodi}{xxxxxxx}


\newif\ifPembimbingHanyaSatu

%%%% WARNING kode berikut ini diaktifkan Jika pembimbing hanya satu %%%%

\PembimbingHanyaSatutrue

%%%%%%%%%%%%%%%%%%%%%%%%%%%%%%%%%%%%%%%%%%%%%%%%%%%%%%%%%%%%%%%%%%

\makeatother
\linespread{1}


\begin{document}
%%%%%%%%%%%%%%%%%%%%%%%%%%%%%%%%%%%%%%%%%%%%%%%%%%%%%%%%%%
% Dimulai dengan cover, lembar persetujuan, Abstrak, dll
%%%%%%%%%%%%%%%%%%%%%%%%%%%%%%%%%%%%%%%%%%%%%%%%%%%%%%%%%%%
\pagenumbering{roman} 
\begin{titlepage}
\thispagestyle{empty}
%\vspace*{0.7cm}
\pdfbookmark{Cover}{ }
{\centering
\large
{\bigsize\bf \Title}\\
\vspace{ 2cm}
\rm
\textbf{Proposal Tugas Akhir}\\
\vspace{0.5 cm}
\textbf{Kelas TA \Kelas}\\
\vspace{0.5 cm}
\textbf{\Author}\\ \textbf{NIM: \NIM}\\ 

\vspace{1.5 cm}

\begin{figure}[h]
{\centering {\includegraphics[scale=0.17]{Tel-U-Logo}}\par}
\end{figure}

\vspace{2 cm}
{\bigsize\textbf{Program Studi Sarjana \Prodi}\\
\vspace{0.5 cm}
\textbf{Fakultas Informatika}\\
\vspace{0.5 cm}
\textbf{Universitas Telkom}\\
\vspace{0.5 cm}
\textbf{Bandung}\\
\vspace{0.5 cm}
\textbf{\Date}\\}
}
\pagebreak
\thispagestyle{empty}
\pdfbookmark{Lembar Persetujuan}{ }
{\centering
\textbf{\large Lembar Persetujuan}\\
\vspace{0.5cm}
\textbf{\Title}\\
\vspace{0.5cm}
\textbf{\textit{\EngTitle}}\\
\vspace{0.5cm}
\textbf{\Author}\\
\textbf{NIM: \NIM}\\
\vspace{1cm}


{ Proposal ini diajukan sebagai usulan pembuatan tugas akhir pada\\ Program Studi Sarjana \Prodi\\ Fakultas Informatika Universitas Telkom}\\

\vspace{0.5cm}

{Bandung, \Tanggal\quad \Bulan \quad \Date}\\
{Menyetujui}\\

\vspace{0.5cm}
      \ifPembimbingHanyaSatu
             \begin{center}
            Calon Pembimbing 1
            \end{center}
       \else
          \begin{center}
              \begin{tabular}{  m{8cm}  m{8cm} }
              Calon Pembimbing 1 & Calon Pembimbing 2
              \end{tabular}
           \end{center}
      \fi
\begin{center}
\vspace{2cm}
\ifPembimbingHanyaSatu
     \underline{\PembimbingSatu} \\ 
      NIP: \NIPPembimbingSatu
\else
\begin{tabular}{  m{8cm}  m{8cm} }
\underline{\PembimbingSatu} & \underline{\PembimbingDua} \\ 
NIP: \NIPPembimbingSatu & NIP: \NIPPembimbingDua
\end{tabular}
\fi
\end{center}
\vspace{0.5cm}

}
\pagebreak
\end{titlepage}
\phantomsection
\addcontentsline{toc}{chapter}{Abstrak}
%\pdfbookmark{\abstractname}{Abstrak Indonesia}
\chapter*{Abstrak}

Riset pemrosesan bahasa natural untuk bahasa Indonesia saat ini terbilang sangat sedikit.
Bahkan, masih banyak area riset yang belum tersentuh seperti contohnya
\textit{combinatory categorial grammar} (CCG).
CCG merupakan formalisme tatabahasa yang pada akhirnya dapat dimanfaatkan untuk memperoleh informasi
dari suatu kalimat.
Informasi tersebut diperoleh setelah melakukan \textit{parsing} berdasarkan formalisme CCG dengan
menggunakan CCG \textit{parser}.
Untuk dapat melakukan \textit{parsing}, CCG \textit{parser} membutuhkan CCG lexicon yang mengandung
bentuk formal dari suatu token kata.
Bentuk formal tersebut adalah \textit{combinatory logic}.
CCG \textit{lexicon} diperoleh dari proses pelabelan suatu token kata terhadap bentuk formalnya dengan
menggunakan \textit{supertagging}.
Proses \textit{supertagging} akan menghasilkan \textit{supertag} yang kemudian disebut sebagai CCG
\textit{lexicon} karena formalisme yang digunakan adalah formalisme CCG.

Tugas akhir dengan judul \textbf{\Title} berusaha untuk membangun versi awal dari CCG
\textit{supertagger} untuk bahasa Indonesia dengan harapan dapat menjadi inisiator untuk riset
pemrosesan bahasa natural dengan tema CCG sehingga ke depannya akan ada lebih banyak riset
mengenai CCG yang tersedia.
\textit{Supertagger} tersebut akan dibangun dengan menggunakan model Maximum Entropy karena
penggunannya di CCG \textit{supertagger} untuk bahasa Inggris sudah memiliki akurasi yang baik.


\vspace{0.5 cm}
\begin{flushleft}
{\textbf{Kata Kunci:}
  combinatory categorial grammar, supertagger, maximum-entropy, bahasa indonesia
}
\end{flushleft}

\cleardoublepage
\phantomsection
\addcontentsline{toc}{chapter}{Daftar Isi}
%\pdfbookmark{\contentsname}{Contents}
\tableofcontents

%%%%%%%%%%%%%%%%%%%%%%%%%%%%%%%%%%%%%%%%%%%%%%%%%
% Bab bab penulisan
%%%%%%%%%%%%%%%%%%%%%%%%%%%%%%%%%%%%%%%%%%%%%%%%
\cleardoublepage
\pagenumbering{arabic}
\chapter{Pendahuluan}
\section{Latar Belakang}

Riset pemrosesan bahasa natural untuk bahasa Indonesia saat ini terbilang sedikit.
Bahkan, masih banyak area riset yang belum tersentuh seperti contohnya
\textit{combinatory categorial grammar} (CCG).
CCG merupakan formalisme tata bahasa yang salah satu manfaatnya adalah untuk memperoleh informasi
(\textit{information extraction}) dari suatu kalimat.
Informasi tersebut diperoleh setelah melakukan \textit{parsing} berdasarkan formalisme CCG dengan
menggunakan perangkat lunak bernama CCG \textit{parser}.
Untuk dapat melakukan \textit{parsing}, CCG \textit{parser} membutuhkan CCG \textit{lexicon}
yang mengandung bentuk formal dari suatu token kata.
Bentuk formal yang dimaksud adalah \textit{category} dalam \textit{category theory}.
CCG \textit{lexicon} diperoleh dari proses pelabelan suatu token kata terhadap bentuk formalnya yang
mana dikenal sebagai \textit{supertagging}.
Proses \textit{supertagging} akan menghasilkan \textit{supertag} yang kemudian disebut sebagai CCG
\textit{supertag} karena formalisme yang digunakan adalah formalisme CCG.
Dalam hal ini, CCG \textit{supertag} adalah CCG \textit{lexicon} itu sendiri.

Tugas akhir dengan judul \say{\Title} berusaha untuk membangun versi awal dari CCG
\textit{supertagger} untuk bahasa Indonesia yang mana harapannya dapat menjadi inisiator riset
pemrosesan bahasa natural dengan tema CCG sehingga ke depannya akan ada lebih banyak riset
mengenai CCG yang tersedia.
\textit{Supertagger} yang dimaksud dalam tugas akhir ini akan dibangun dengan menggunakan model
Maximum Entropy (MaxEnt) dan implementasinya akan ditulis dalam bahasa pemrograman Haskell.
Model MaxEnt digunakan karena keterbatasan \textit{dataset} untuk melakukan \textit{learning}.
Adapun bahasa pemrograman Haskell digunakan karena abstraksi bahasanya yang sangat mendekati
\textit{category theory} serta kemampuannya yang sangat baik dalam pemrosesan data.

\section{Perumusan Masalah}
Rumusan masalah yang akan diangkat yaitu:
\begin{enumerate}
    \item Mengapa CCG \textit{supertagger} diperlukan?
    \item Apa saja yang harus dipersiapkan untuk membangun CCG \textit{supertagger}?
    \item Bagaimana proses pembangunan CCG \textit{supertagger}?
\end{enumerate}
\section{Tujuan}
Tujuan yang diharapkan dapat tercapai oleh tugas akhir ini yaitu:
\begin{enumerate}
    \item Mengenalkan alternatif metode yang dapat digunakan dalam pemrosesan bahasa alami untuk
      bahasa Indonesia.
    \item Merilis CCG \textit{supertagger} pertama untuk bahasa Indonesia.
    \item Membuka peluang riset untuk CCG \textit{parser} bahasa Indonesia.
\end{enumerate}
\section{Batasan Masalah}
Hipotesis dari tugas akhir ini yaitu:
\begin{enumerate}
    \item Memberikan label CCG untuk proses \textit{learning} merupakan permasalahan utama dari tugas
      akhir ini.
    \item \textit{Supertagger} yang akan dibangun kemungkinan besar memiliki akurasi yang cenderung
      rendah.
    \item CCG \textit{lexicon} sudah dapat digunakan oleh CCG \textit{parser} bahasa Indonesia
      (apabila ada).
\end{enumerate}

\section{Rencana Kegiatan}
Rencana kegiatan yang akan dilakukan adalah sebagai berikut:
\begin{itemize}
    \item Studi literatur
    \item Studi \textit{tools} yang tersedia
    \item Studi bahasa pemrograman yang akan digunakan
    \item Perancangan sistem \textit{supertagger}
    \item Membangun \textit{supertagger}
    \item Memeriksa hasil
\end{itemize}

\section{Jadwal Kegiatan}
Laporan proposal ini akan dijadwalkan sesuai dengan tabel \ref{ScheduleTable}. 
\begin{table}[h!]
  \centering
    \caption{Jadwal kegiatan proposal tugas akhir.}
  \label{ScheduleTable}
  \begin{tabular}{|c|m{2.5cm}|m{0.01cm}|m{0.01cm}|m{0.01cm}|m{0.01cm}|m{0.01cm}|m{0.01cm}|m{0.01cm}|m{0.01cm}|m{0.01cm}|m{0.01cm}|m{0.01cm}|m{0.01cm}|m{0.01cm}|m{0.01cm}|m{0.01cm}|m{0.01cm}|m{0.01cm}|m{0.01cm}|m{0.01cm}|m{0.01cm}|m{0.01cm}|m{0.01cm}|m{0.01cm}|m{0.01cm}|}
    \hline
    \multirow{2}{*}{\textbf{No}} & \multirow{2}{*}{\textbf{Kegiatan}} & \multicolumn{24}{|c|}{\textbf{Bulan ke-}} \\
    \hhline{~~------------------------}
    {} & {} & \multicolumn{4}{|c|}{\textbf{1}} & \multicolumn{4}{|c|}{\textbf{2}} & \multicolumn{4}{|c|}{\textbf{3}} & \multicolumn{4}{|c|}{\textbf{4}} & \multicolumn{4}{|c|}{\textbf{5}} & \multicolumn{4}{|c|}{\textbf{6}}\\
    \hline
    1 & Studi Literatur & \cellcolor{blue!25} & \cellcolor{blue!25} & \cellcolor{blue!25} & \cellcolor{blue!25}& \cellcolor{blue!25} & \cellcolor{blue!25} & \cellcolor{blue!25} & \cellcolor{blue!25}& \cellcolor{blue!25} & \cellcolor{blue!25} & \cellcolor{blue!25} & \cellcolor{blue!25}& \cellcolor{blue!25} & \cellcolor{blue!25} & \cellcolor{blue!25} & \cellcolor{blue!25}& \cellcolor{blue!25} & \cellcolor{blue!25} & \cellcolor{blue!25} & \cellcolor{blue!25}& \cellcolor{blue!25} & \cellcolor{blue!25} & \cellcolor{blue!25} & \cellcolor{blue!25}\\
    \hline
    2 & Studi \textit{Tools} yang Tersedia & \cellcolor{blue!25} & \cellcolor{blue!25} & \cellcolor{blue!25} & \cellcolor{blue!25} &  \cellcolor{blue!25} &  \cellcolor{blue!25} &  \cellcolor{blue!25} &  \cellcolor{blue!25} & {} & {} & {} & {}& {} & {} & {} & {}& {} & {} & {} & {}& {} & {} & {} & {}\\
    \hline
    3 & Studi Bahasa Pemrograman & {} & {} & {} & {} & \cellcolor{blue!25} & \cellcolor{blue!25} & \cellcolor{blue!25} & \cellcolor{blue!25} & {} & {} & {} & {}& {} & {} & {} & {}& {} & {} & {} & {}& {} & {} & {} & {}\\
    \hline
    4 & Pengumpulan Data & {} & {} & {} & {} & \cellcolor{blue!25} & \cellcolor{blue!25} & \cellcolor{blue!25} & \cellcolor{blue!25} & \cellcolor{blue!25} & \cellcolor{blue!25} & \cellcolor{blue!25} & \cellcolor{blue!25} & {} & {} & {} & {}& {} & {} & {} & {} & {} & {} & {} & {}\\
    \hline
    5 & Analisis dan Perancangan Sistem &  {} & {} & {} & {}  & \cellcolor{blue!25} & \cellcolor{blue!25} & \cellcolor{blue!25} & \cellcolor{blue!25} & \cellcolor{blue!25} & \cellcolor{blue!25} & \cellcolor{blue!25} & \cellcolor{blue!25} & {} & {} & {} & {}& {} & {} & {} & {}& {} & {} & {} & {}\\
    \hline
    6 & Implementasi Sistem &  {} & {} & {} & {} & {} & {} & {} & {}& \cellcolor{blue!25} & \cellcolor{blue!25} & \cellcolor{blue!25} & \cellcolor{blue!25} & \cellcolor{blue!25} & \cellcolor{blue!25} & \cellcolor{blue!25} & \cellcolor{blue!25} & {} & {} & {} & {}& {} & {} & {} & {}\\
    \hline
    7 & Analisa Hasil Implementasi &  {} & {} & {} & {} & {} & {} & {} & {}& {} & {} & {} & {} & \cellcolor{blue!25} & \cellcolor{blue!25} & \cellcolor{blue!25} & \cellcolor{blue!25} & \cellcolor{blue!25} & \cellcolor{blue!25} & \cellcolor{blue!25} & \cellcolor{blue!25} & {} & {} & {} & {}\\
    \hline
    8 & Penulisan Laporan & {} & {} & {} & {} & \cellcolor{blue!25} & \cellcolor{blue!25} & \cellcolor{blue!25} & \cellcolor{blue!25}& \cellcolor{blue!25} & \cellcolor{blue!25} & \cellcolor{blue!25} & \cellcolor{blue!25}& \cellcolor{blue!25} & \cellcolor{blue!25} & \cellcolor{blue!25} & \cellcolor{blue!25}& \cellcolor{blue!25} & \cellcolor{blue!25} & \cellcolor{blue!25} & \cellcolor{blue!25}& \cellcolor{blue!25} & \cellcolor{blue!25} & \cellcolor{blue!25} & \cellcolor{blue!25}\\
    \hline
  \end{tabular}
\end{table}

%
\chapter{Kajian Pustaka}

\section{Categorial Grammar}
Categorial Grammar (CG) merupakan sebuah istilah yang mencakup beberapa formalisme terkait yang diajukan
untuk sintaks dan semantik dari bahasa alami serta untuk bahasa logis dan matematis \cite{Steedman92catg}.
Karakteristik yang paling terlihat dari CG adalah bentuk esktrim dari leksikalismenya di mana beban utama
(atau bahkan seluruh beban) sintaksisnya ditanggung oleh leksikon.
Konstituen tata bahasa dalam \textit{categorial grammar} dan khususnya semua leksikal diasosiasikan
dengan suatu \textit{type} atau \say{\textit{category}} (dalam \textit{category theory}) yang
mendefinisikan potensi mereka untuk dikombinasikan dengan konstituen lain untuk menghasilkan konstituen
majemuk.
\textit{Category} tersebut adalah salah satu dari sejumlah kecil \textit{category} dasar (seperti NP)
atau \textit{functor} (dalam \textit{category theory}).

Ada beberapa notasi berbeda untuk \textit{category} dalam merepresentasikan \textit{directional}-nya.
Notasi yang paling umum digunakan adalah \say{\textit{slash notation}} yang dipelopori oleh Bar-Hilel,
Lambek, dan kemudian dimodifikasi dalam kelompok teori yang dibedakan sebagai tata bahasa
\say{\textit{combinatory}} \textit{categorial grammar} (CCG).
Sebagai contoh, \textit{category} $\text{(S$\backslash$NP)/NP}$ merupakan suatu \textit{functor} yang
memiliki dua buah notasi \textit{slash} yaitu $\backslash$ dan $/$.
Masing-masing notasi \textit{slash} tersebut merepresentasikan \textit{directionality} yang berbeda.
Notasi \textit{forward slash}, $/$, mengindikasikan bahwa argumen dari suatu \textit{functor}
$\text{X}/\text{Y}$ ada di bagian kanan atau dengan kata lain $\text{Y}$.
Adapun \textit{backward slash}, $\backslash$, mengindikasikan bahwa argumen dari suatu \textit{functor}
$\text{X}\backslash\text{Y}$ ada di bagian kiri atau dengan kata lain $\text{X}$.
Demikian itu, penggunaan notasi \textit{slash} yang tepat sangat penting dikarenakan hal ini dapat
mempengaruhi konstituen dari hasil \say{kombinasi} \textit{category}-nya.

\section{Combinatory Categorial Grammar}
Combinatory Categorial Grammar (CCG) merupakan salah satu formalisme tata bahasa yang gaya aturannya
diturunkan dari \textit{categorial grammar} dengan beberapa penambahan aturan dan istilah baru.
Di CCG, \textit{category} dapat dipasangkan dengan \textit{semantic representation}.
Dalam hal ini, \textit{semantic representation} yang dimaksud adalah abstraksi fungsi lambda
(dalam \textit{lambda calculus}, \textit{lambda function}).
Sebagai contoh, \textit{category} $\text{(S$\backslash$NP)/NP}$ dapat dipasangkan dengan fungsi lambda
$\lambda{x. fx}$ sehingga dapat ditulis menjadi $\text{(S$\backslash$NP)/NP} : \lambda{x. fx}$.
Adapun pemetaan dari suatu token kata ke \textit{category}-nya menggunakan notasi $\vdash$.
Sebagai contoh, anggap saja kita memiliki kamus pemetaan seperti pada Gambar \ref{ccg:mapping:1}.
Apabila kita memiliki kalimat \say{Pamungkas dan Setyo menyukai rendang}, maka kita dapatkan:

\begin{figure}\centering\small
  \begin{align*}
    \text{Pamungkas} &\vdash \text{NP}: \so{pamungkas}\\
    \text{Setyo} &\vdash \text{NP}: \so{setyo}\\
    \text{dan} &\vdash \text{CONJ}: \lambda x.\lambda y.\lambda f.\ (f\ x) \land (f\ y)\\
    \text{menyukai} &\vdash \text{(S{$\backslash$}NP)/NP}: \lambda x.\lambda y.\ suka(y, x)\\
    \text{rendang} &\vdash \text{NP}: \so{rendang}
  \end{align*}
  \caption{Kamus yang memetakan token kata ke bentuk CCG \textit{lexicon}-nya.}
  \label{ccg:mapping:1}
\end{figure}

\begin{center}
  \bgroup
  \catcode`!=\active \def!{\upshape}
  \catcode`?=\active \def?#1{\makebox[0pt]{#1}}
  \catcode`^=\active \def^#1{\footnotesize{#1}}
  \catcode`*=\active \def*#1{\scriptsize{#1}}
  \tabbedShortstack{
    !^Pamungkas & & !^dan & & !^Setyo & & !^menyukai & & !^rendang &\\
    \TABcline{1,3,5,7,9}
    !^{$\text{NP}$} & &
      !^{$\text{CONJ}$} & &
      !^{$\text{NP}$} & &
      !^{$\text{(S$\backslash$NP)/NP}$} & &
      !^{$\text{NP}$} &\\
    !{*: \so{pamungkas}} & &
      !{*: $\lambda x.\lambda y.\lambda f.\ (f\ x) \land (f\ y)$} & &
      !{*: \so{setyo}} & &
      !{*: $\lambda x.\lambda y.\ suka(y, x)$} & &
      !{*: \so{rendang}} &
  }
  \egroup
\end{center}

Ada beberapa operasi yang dapat dilakukan dalam CCG. \textit{Operand} dari operasi
yang dimaksud adalah \textit{category}. Berdasarkan contoh di atas, akan ada tiga
operasi yang dijalankan yaitu \textit{coordination}, \textit{forward application},
dan \textit{type rising}.
Untuk mendapatkan hasil yang diinginkan, kita lakukan \textit{type rising} sebelum
\textit{forward application} di akhir.
Sehingga, kita dapatkan:

\begin{center}
  \bgroup
  \catcode`!=\active \def!{\upshape}
  \catcode`?=\active \def?#1{\makebox[0pt]{#1}}
  \catcode`^=\active \def^#1{\footnotesize{#1}}
  \catcode`*=\active \def*#1{\scriptsize{#1}}
  \tabbedLongunderstack{
    !^Pamungkas & & !^dan & & !^Setyo & & !^menyukai & & !^rendang &\\
    \TABcline{1,3,5,7,9}
    !^{$\text{NP}$} & &
      !^{$\text{CONJ}$} & &
      !^{$\text{NP}$} & &
      !^{$\text{(S$\backslash$NP)/NP}$} & &
      !^{$\text{NP}$} &\\
    !{*: \so{pamungkas}} & &
      !{*: $\lambda x.\lambda y.\lambda f.\ (f\ x) \land (f\ y)$} & &
      !{*: \so{setyo}} & &
      !{*: $\lambda x.\lambda y.\ suka(y, x)$} & &
      !{*: \so{rendang}} &\\
    \TABrule & \TABrule &
      \TABrule & \TABrule &
      \TABrule\CCGCOOR & &
      \TABrule & \TABrule &
      \TABrule\CCGFA &\\
    & &
      !?{^{$\text{NP}$}}
      \ \ \ \ \ \ \ \ \ 
      & & & &
      \ \ \ \ \ \ \ 
      !?{^{$\text{S$\backslash$NP}$}}
      & & &\\
    & &
      ?{*: $\lambda f.\ (f\ \so{pamungkas}) \land (f\ \so{setyo})$}
      \ \ \ \ \ \ \ \ \ 
      & & & &
      \ \ \ \ \ \ \ 
      ?{*: $\lambda y.\ suka(y, \so{rendang})$}
      & & &\\
    \TABrule & \TABrule &
      \TABrule & \TABrule &
      \TABrule\CCGTR & &
      & & &\\
    & &
      !?{^{$\text{S/(S$\backslash$NP)}$}}
      \ \ \ \ \ \ \ \ \ 
      & & & & & & &\\
    & &
      ?{*: $\lambda f.\ (f\ \so{pamungkas}) \land (f\ \so{setyo})$}
      \ \ \ \ \ \ \ \ \ 
      & & & & & & &\\
    \TABrule & \TABrule &
      \TABrule & \TABrule &
      \TABrule & \TABrule &
      \TABrule & \TABrule &
      \TABrule\CCGFA &\\
    & & &
      !?{^{$\text{S}$}}
      & & & & & &\\
    & & &
      ?{*: $suka(\so{pamungkas}, \so{rendang}) \land suka(\so{setyo}, \so{rendang})$}
      & & & & & &
  }
  \egroup
\end{center}

Berdasarkan hasil evaluasi tersebut, kita dapatkan \textit{query} \ref{ccg:query:1}
yang diperoleh dari kalimat \say{Pamungkas dan Setyo menyukai rendang}.
Demikian itu, komputer dapat melakukan komputasi berdasarkan \textit{query} yang telah diperoleh.
Kegiatan tersebut merupakan apa yang disebut dengan CCG \textit{parsing}.
Untuk dapat melakukan parsing, CCG \textit{lexicon} diperlukan.
Untuk mendapatkan CCG \textit{lexicon} kita dapat menggunakan CCG \textit{supertagger}
yang akan melakukan pelabelan suatu token kata ke CCG \textit{lexicon} berdasarkan
pemetaannya.

\begin{equation}\label{ccg:query:1}
  suka(\so{pamungkas}, \so{rendang}) \land suka(\so{setyo}, \so{rendang})
\end{equation}

\section{Category Theory}
\textit{Category Theory} (CT) merupakan formalisme yang dapat digunakan untuk memformalkan
struktur matematis.
CT mempelajari \textit{category} yang merupakan sebuah representasi dari suatu
abstraksi konsep matematis.
Suatu \textit{category} memiliki kumpulan \textit{object} dan \textit{morphism}.
Untuk mempermudah pemahaman mengenai CT, kita akan gunakan
\textit{category of set} (kategori dari himpunan) sebagai contoh.
Dalam \textit{category of set}, \textit{object}-nya adalah himpunan dan
\textit{morphism}-nya (terkadang disebut dengan \textit{arrow}) adalah fungsi
(\textit{function}, sebuah pemetaan).
Kemudian, pemetaan dari suatu \textit{category} $C$ ke \textit{category} $D$
yang dipetakan oleh $F$ ($F: C \rightarrow D$) disebut sebagai \textit{functor}.


\section{Lambda Calculus}
\textit{Lambda calculus} ({$\lambda$}\textit{-calculus}) merupakan sebuah formalisme yang dikembangkan
oleh Alonzo Church sebagai alat yang digunakan untuk memahami konsep komputasi yang efektif
\cite{DBLP:journals/corr/Rojas15}.
Formalisme {$\lambda$}\textit{-calculus} cukup populer dan bahkan dijadikan sebagai pondasi teori bagi
paradigma pemrograman \textit{functional programming}.
Konsep utama dari {$\lambda$}\textit{-calculus} adalah apa yang disebut dengan \textit{expression}.
Suatu \textit{expression} dalam {$\lambda$}\textit{-calculus} terdiri dari tiga bagian yaitu
\textit{lambda notation} ({$\lambda$}), \textit{argument} (seperti $a$, $b$, $c$, $x$, dan lain-lain),
dan \textit{body} yang dipisahkan dengan tanda titik.
Sebagai contoh, fungsi lambda ${\lambda}x. x$ merupakan sebuah fungsi identitas yang mengambil
argumen $x$ kemudian mengembalikan nilai $x$ itu sendiri.
Dalam hal ini, terlihat bahwa notasi {$\lambda$} merupakan sebuah penanda bagi suatu fungsi lambda.
Kemudian, pengubah $x$ setelah notasi {$\lambda$} merupakan argumen dari fungsi tersebut.
Selanjutnya, tanda titik merupakan pemisah antara \textit{head} dan \textit{body} fungsi lambda.
Terakhir, setelah tanda titik adalah \textit{body} dari suatu fungsi lambda yang mana berupa
\textit{expression}.

Untuk mempermudah pemahaman, {$\lambda$}\textit{-calculus} dapat diperlakukan seperti fungsi tanpa
nama. Sebagai contoh, fungsi lambda $({\lambda}x. x + 5)$ apabila diberikan nilai $2$ sehingga
menjadi $({\lambda}x. x + 5) 2$ akan dievaluasi menjadi ${\lambda}(2). (2) + 5$.
Demikian itu, nilai yang dikembalikan oleh fungsi tersebut adalah $7$.
Sama seperti fungsi pada umumnya, konsep ini bernama \textit{substition} (substitusi).
Memahami {$\lambda$}\textit{-calculus} dirasa perlu berhubung dalam tugas akhir ini
{$\lambda$}\textit{-calculus} digunakan sebagai bentuk formal di \textit{category}
dalam konteks CCG \textit{lexicon}. Meskipun {$\lambda$}\textit{-calculus} tidak sesederhana
yang dijelaskan sebelumnya, setidaknya memahami {$\lambda$}\textit{-calculus} seperti ini
sudah cukup untuk dapat membangun \textit{supertagger} yang ada di tugas akhir ini.


\section{Supertagging}
\textit{Supertagging} merupakan proses yang memetakan suatu token kata ke bentuk
\textit{supertag}-nya. CCG \textit{supertagging} artinya proses pemetaan dari suatu token
kata ke dalam bentuk CCG \textit{supertag} atau dapat juga disebut sebagai CCG \textit{lexicon}.
\textit{Supertag} mirip dengan POS \textit{tag}.
Perbedaannya, \textit{supertag} memiliki bentuk formalisme yang lebih kompleks dari POS
\textit{tag}.
Hal ini dikarenakan \textit{supertag} menyimpan informasi lain selain tanda gramatikalnya
($\text{NP}$, $\text{NN}$, $\text{VB}$, dan sebagainya).
Sebagai contoh, CCG \textit{supertag} memiliki bentuk dengan format:

\begin{equation*}
<\text{categorial\ grammar\ tag}>:\ <\text{semantic\ representation}>
\end{equation*}

Kita dapat menggunakan notasi $\vdash$ untuk memetakan token kata ke bentuk CCG
\textit{supertag}-nya.
Anggap kita memiliki kata kerja \say{menyukai} yang akan dipetakan ke
$\text{(S{$\backslash$}NP)/NP}: \lambda x.\lambda y.\ suka(y, x)$, kita dapatkan:

\begin{equation*}
  \text{menyukai} \vdash (S{\backslash}NP)/NP: \lambda x.\lambda y.\ suka(y, x)
\end{equation*}

Adapun sebuah \textit{tool} yang melakukan proses \textit{supertagging} ini dinamakan
\textit{supertagger}.
\textit{Supertagger} sederhananya mengambil daftar token kata yang kemudian untuk setiap
token kata tersebut \say{dilabelkan} dengan \textit{supertag}-nya.


\section{Maximum Entropy Model}
\textit{Maximum Entropy} (MaxEnt) merupakan model statistik yang dapat digunakan untuk
melakukan \textit{train} korpus teranotasi dengan Part-Of-Speech (POS) \textit{tag}.
Salah satu aplikasi MaxEnt adalah POS \textit{tagger} yang mana memiliki hasil akurasi
yang lebih baik ketimbang \textit{state-of-the-art}.
Adapun model probabilitasnya didefinisikan dalam $\mathcal{H}\times\mathcal{T}$,
dimana $\mathcal{H}$ adalah himpunan dari kemungkinan kata dan konteks \textit{tag},
atau \say{\textit{history}} (riwayat), dan $\mathcal{T}$ adalah himpunan dari
\textit{tag} yang diizinkan.
Model probabilitas dari suatu \textit{history} $h$ bersama dengan \textit{tag} $t$
didefinisikan dalam persamaan \ref{maxent:equation:1}.

\begin{equation}\label{maxent:equation:1}
  p(h, t) = \pi\mu\prod_{j = 1}^{k} a_{j}^{f_{j}(h, t)}
\end{equation}

\noindent
dimana $\pi$ merupakan konstan normalisasi,
$\{\mu, a_1, \dots, a_k\}$ merupakan parameter model positif,
dan $\{f_1, \dots, f_k\}$ merupakan apa yang kita sebut sebagai
\say{\textit{feature}},
dimana $f_j(h, t) \in \{0, 1\}$.

Meskipun MaxEnt dipernalkan untuk POS \textit{tagging},
MaxEnt dapat pula digunakan untuk \textit{supertagging}.
Stephen Clark 2002 mempublikasikan literatur \textit{supertagging} untuk CCG yang
mana MaxEnt merupakan model yang digunakan.
Adapun persamaan modelnya dapat dilihat di persamaan \ref{maxent:equation:2}.

\begin{equation}\label{maxent:equation:2}
  p(c|h) = \frac{1}{Z(h)} e^{\Sigma_i \lambda_if_i(c, h)}
\end{equation}

\noindent
dimana $c$ merupakan \textit{category}, $h$ merupakan \textit{context},
fungsi $f_i(c, h)$ merupakan \say{\textit{feature}} dari suatu
\textit{category} dan \textit{context}, dan $Z(h)$ merupakan konstan normalisasinya.
Adapun contoh \say{\textit{feature}} yang dimaksud dapat dilihat di persamaan
\ref{maxent:equation:3}.

\begin{equation}\label{maxent:equation:3}
  f_j(c, h) =
  \begin{cases}
    1 &, \text{if } merupakan\_kata\_yang(h) = true\ \&\ c = \text{NP/N}\\
    0 &, \text{selainnya}
  \end{cases}
\end{equation}


% \section{Time Seies method}
% Menurut paper Kentang \cite{Kentang}, prsamaan SWE adalah
% Berikut diberikan persamaan pengatur dari persamaan gelombang pada gitar

% \begin{equation}\label{Pers1}
%     a=b+U^{n+1}_{i+1}
% \end{equation}
% Persamaan (\ref{Pers1}) jadsbahdhavhdvah ajdbajdb
% \begin{equation}\label{nama-rumus}
%     \int_0^1 \frac{f(x)}{g(x)}\ {\rm dx}=\sin x
% \end{equation}

% \begin{equation}\label{nama-rumus1}
%    \alpha \times \beta =\gamma^{3\alpha}
% \end{equation}

% \begin{figure}[h!]
%     \centering
%     \includegraphics[scale=0.1]{Tel-U-Logo.png}
%     \caption{Caption}
%     \label{fig:my_label1}
% \end{figure} 

% Rumus (\ref{nama-rumus}) merupakan \textit{contoh} persamaan matematika. persamaan matematika diatas diberi nama \textbackslash label\{nama-rumus\}. dengan $\alpha=\gamma \times 100$

% \begin{figure}[h!]
%     \centering
%     \includegraphics[scale=0.3]{Tel-U-Logo.png}
%     \caption{Caption}
%     \label{fig:my_label}
% \end{figure}

% Lihat \textit{pada} Gambar \ref{fig:my_label}

% \subsection{Cara memanggil pustaka}
% Contoh pustaka prosiding \cite{doyen2014explicit}, jurnal \cite{gunawan2015hydrostatic} dan buku \cite{toro2013riemann}. Atau dapat juga mengguanakan dua pustaka atau lebih dalam \cite{gunawan2015hydrostatic,toro2013riemann}.
%
% \tikzstyle{startstop} = [rectangle, rounded corners, minimum width=3cm, minimum height=1cm,text centered, draw=black, fill=red!30]
\tikzstyle{io} = [trapezium, trapezium left angle=70, trapezium right angle=110, minimum width=3cm, minimum height=1cm, text centered, draw=black, fill=blue!30]
\tikzstyle{process} = [rectangle, minimum width=3cm, minimum height=1cm, text centered, minimum width=3cm, draw=black, fill=orange!30]
\tikzstyle{decision} = [diamond, minimum width=3cm, minimum height=1cm, text centered, draw=black, fill=green!30]
\tikzstyle{arrow} = [thick,->,>=stealth]
\tikzstyle{cloud} = [draw, ellipse,fill=red!20, minimum height=2em]

\chapter{Metodologi dan Desain Sistem}

\section{Pembuatan Dataset}

\textit{Dataset} untuk melakukan \textit{train} CCG sayangnya belum tersedia.
Karenanya, kita akan membuatkan \textit{dataset} secara semi otomatis dengan memanfaatkan POS
\textit{tagger} untuk bahasa Indonesia.
Secara formal, \textit{flowchart} untuk proses pembuatan \textit{dataset} dapat dilihat pada
Gambar \ref{flowchart:dataset:1}.

\begin{figure}[h!]\centering
	\begin{tikzpicture}[node distance=2cm]
		\node (start) [cloud] {Start};
		\node (in) [io, below of=start, yshift=0.5cm] {\textit{Input} teks bahasa Indonesia};
		\node (pro1) [process, below of=in, yshift=0.5cm] {POS \textit{tagging}};
		\node (pro2) [process, below of=pro1, yshift=0.5cm] {Jadikan \textit{queue}};
		\node (dec) [decision, below of=pro2, yshift=-1cm] {\textit{Queue} kosong?};
		\node (pro3) [process, below of=dec, yshift=-1cm] {Transformasi POS \textit{tag}};
		\node (out) [io, below of=pro3, yshift=0.5cm] {\textit{Output} CCG \textit{supertag}};
		\node (stop) [cloud, below of=out, yshift=0.5cm] {Stop};
		\draw [arrow] (start) -- (in);
		\draw [arrow] (in) -- (pro1);
		\draw [arrow] (pro1) -- (pro2);
		\draw [arrow] (pro2) -- (dec);
		\draw [arrow] (dec) -- node[anchor=east]{Tidak} (pro3);
		\draw [arrow] (pro3) -- +(-3,0) |- (dec);
		\draw [arrow] (dec) -- node[anchor=south]{Ya} +(4,0) |- (out);
		\draw [arrow] (out) -- (stop);
	\end{tikzpicture}
	\caption{Alur kerja pembuatan \textit{dataset} untuk CCG \textit{supertag}.}
	\label{flowchart:dataset:1}
\end{figure}

\subsection{Input Teks Bahasa Indonesia}

Pada bagian \textit{input} dalam pembuatan \textit{dataset}, kita dapat mengambil teks bahasa Indonesia
dari Indonesian Treebank\footnotemark[1]\ dan/atau dari beberapa contoh artikel yang terdapat di
\textit{website} Wikipedia Indonesia\footnotemark[2].
Indonesian Treebank terdapat setidaknya 1000 kalimat yang dirasa cukup untuk melakukan
\textit{training} menggunakan model MaxEnt.
Apa yang perlu dilakukan setelah mengambil \textit{input} adalah membersihkan bentuk \textit{tree}-nya
untuk kemudian diambilkan kalimat yang sebenarnya.
Caranya cukup sederhana yaitu dengan mengambil \textit{leaf} dari \textit{tree} tersebut kemudian
disatukan di suatu pengubah dengan tipe data \textit{string}.

Sebagai contoh, salah satu \textit{tree} yang terdapat di Indonesian Treebank dapat dilihat pada
Gambar \ref{treebank:tree:1}.
Apa yang dimaksud dengan \textit{leaf} pada \textit{tree} tersebut adalah $(Kera)$, $(untuk)$,
$(*)$, $(amankan)$, $(pesta\ olahraga)$.
Terkhusus untuk \textit{leaf} dengan bentuk spesial, seperti $(*)$, kita hilangkan sehingga teks
yang diperoleh dari \textit{tree} tersebut adalah \say{Kera untuk amankan pesta olahraga}.
Selanjutnya, contoh teks yang akan digunakan agar konsisten yaitu
\say{Pamungkas dan Setyo menyukai rendang}.


\begin{figure}\centering\small
	\begin{align*}
		&(NP\\
		&\ \ (NN\ (Kera))\\
		&\ \ (SBAR\\
		&\ \ \ \ (SC\ (untuk))\\
		&\ \ \ \ (S\ (NP-SBJ\ (*))\\
		&\ \ \ \ \ \ (VP\\
		&\ \ \ \ \ \ \ \ (VB\ (amankan))\\
		&\ \ \ \ \ \ \ \ (NP\ (NN\ (pesta\ olahraga)))))))
	\end{align*}
	\caption{Salah satu contoh \textit{tree} dalam Indonesian Treebank.}
	\label{treebank:tree:1}
\end{figure}

\footnotetext[1]{github.com/famrashel/idn-treebank/blob/master/Indonesian\_Treebank.bracket}
\footnotetext[2]{id.wikipedia.org}

\subsection{POS Tagging}

Berdasarkan teks bahasa Indonesia yang telah diambil, kita manfaatkan \textit{tool} POS \textit{tagger}
bahasa Indonesia untuk mendapatkan \textit{lexical category} atomik untuk masing-masing token.
Dengan memanfaatkan POS \textit{tag} kita dapat membuat \textit{dataset} untuk CCG \textit{supertag}
lebih mudah dibandingkan dengan memberikan \textit{tag} CCG secara manual.
Ide dasarnya yaitu kita akan mentransformasikan POS \textit{tag} yang didapatkan menjadi CCG
\textit{supertag} berdasarkan aturan-aturan khusus yang telah ditentukan.
Sebagai contoh, kita dapat membuat aturan seperti mentransformasikan $\text{VB}$ menjadi
$\text{(S$\backslash$NP)/NP}$.
Dengan menggunakan contoh kalimat yang sama seperti di bagian sebelumnya,
yaitu \say{Pamungkas dan Setyo menyukai rendang}, setelah menggunakan POS \textit{tagger} bahasa
Indonesia kita dapatkan hasil sesuai dengan Gambar \ref{postag:1}.

\begin{figure}\centering
  \bgroup
  \catcode`!=\active \def!{\upshape}
  \catcode`?=\active \def?#1{\makebox[0pt]{#1}}
  \catcode`^=\active \def^#1{\footnotesize{#1}}
  \catcode`*=\active \def*#1{\scriptsize{#1}}
  \tabbedShortstack{
    !^Pamungkas & & !^dan & & !^Setyo & & !^menyukai & & !^rendang &\\
    \TABcline{1,3,5,7,9}
    !^{$\text{NNP}$} & &
      !^{$\text{CC}$} & &
      !^{$\text{NNP}$} & &
      !^{$\text{VB}$} & &
      !^{$\text{X}$} &
  }
	\egroup
	\caption{Kalimat contoh dengan POS \textit{tag}-nya.}
	\label{postag:1}
\end{figure}

\subsection{Transformasi POS Tag}

Pada bagian proses transformasi POS \textit{tag} ke bentuk CCG \textit{supertag}-nya,
kita buatkan aturan-aturan transformasinya.
Aturan transformasi tersebut merupakan pemetaan berbasis aturan.
Sebagai contoh, kata \say{menyukai} memiliki POS \textit{tag} $\text{VB}$.
Anggap saja dalam aturan transformasi terdapat pemetaan $\text{VB} \vdash \text{(S$\backslash$NP)/NP}$.
Sehingga, kita dapatkan CCG \textit{supertag} untuk \say{menyukai} yaitu $\text{(S$\backslash$NP)/NP}$.
Kendati demikian, masih ada bagian yang belum kita dapatkan yaitu \textit{semantic representation}-nya.
Kita dapat menggunakan \textit{stemmer} bahasa Indonesia agar mendapatkan \textit{root words}
dari kata \say{menyukai} yaitu \say{suka}.
Langkah terakhirnya adalah membuatkan \textit{semantic representation}-nya berdasarkan
\textit{root words} yang telah diperoleh sehingga kita dapatkan fungsi lambdanya yaitu
$\lambda x.\lambda y.\ \text{suka}(y, x)$.

Selain menggunakan \textit{stemmer}, kita dapat menggunakan \textit{morphological analyzer}.
Untuk bahasa Indonesia, kita dapat menggunakan \textit{tool} bernama MorphInd\footnotemark[3].
\textit{Morphological analyzer} salah satu kegunaannya yaitu dapat menghasilkan \textit{root words}
sehingga dapat kita manfaatkan untuk membuat fungsi lambda.
Namun, kita dapat menggunakan MorphInd sebagai pelengkap POS \textit{tagger} untuk bahasa Indonesia.
Hal ini agar \textit{dataset} yang dibuatkan secara semi-otomatis ini dapat memiliki kualitas yang
baik.

\footnotetext[3]{septinalarasati.com/morphind}

\subsection{Output CCG Supertag}

Keluaran dari bagian pembuatan \textit{dataset} ini adalah sebuah berkas berisi CCG \textit{supertag}
lengkap dan beberapa berkas CCG \textit{supertag} yang lemanya dikelompokkan berasarkan alfabet.
Adapun format \textit{dataset} yang disimpan dalam berkas tersebut dapat dilihat pada Gambar
\ref{dataset:format}.
Untuk setiap kalimat dalam berkas tersebut direpresentasikan oleh dua baris yaitu
(1) daftar token yang dipisahkan oleh spasi, dan
(2) daftar \textit{supertag} yang dibungkus kurung siku (contoh: $\text{[\textit{supertag}]}$) dan
dipisahkan oleh spasi sejajar dengan tokennya.
Sebagai contoh, dengan kalimat yang sama dengan sebelum-sebelumnya, dapat dilihat di Gambar
\ref{dataset:example}.
Demikian itu, kita dapat mengambil \textit{dataset} secara lengkap dengan CCG \textit{supertag}-nya.

\begin{figure}\centering
	\bgroup
	\tabbedShortstack{
		$\text{token}_1$ & $\text{token}_2$ & $\text{token}_3$ & $\dots$ & $\text{token}_n$\\
		$\text{supertag}_1$ &
			$\text{supertag}_2$ & $\text{supertag}_3$ & $\dots$ & $\text{supertag}_n$
	}
	\egroup
	\caption{Format \textit{dataset} yang disimpan di dalam berkas.}
	\label{dataset:format}
\end{figure}

\begin{figure}\centering\tiny
	\bgroup
	\tabbedShortstack{
		Pamungkas & dan & Setyo & menyukai & rendang & .\\
		$\text{[NP\ :\ \so{pamungkas}]}$ &
			$\text{[CONJ\ :\ $\lambda x.\lambda y.\lambda f.\ (f\ x) \land (f\ y)$]}$ &
			$\text{[NP\ :\ \so{Setyo}]}$ &
			$\text{[(S{$\backslash$}NP)/NP\ :\ $\lambda x.\lambda y.\ suka(y, x)$]}$ &
			$\text{[NP\ :\ \so{rendang}]}$ &
			$\text{[Z]}$
	}
	\egroup
	\caption{Contoh isi dari berkas \textit{dataset}.}
	\label{dataset:example}
\end{figure}

\section{Melatih Model Supertagger}

Model yang akan digunakan oleh \textit{supertagger} ini adalah Maximum Entropy (MaxEnt).
Model ini dipilih karena ketersediaan \textit{dataset} bahasa Indonesia yang masih sangat kurang.
Selain itu, model MaxEnt digunakan di sebuah riset yang dilakukan oleh Stephen Clark dalam
pembuatan \textit{supertagger}-nya.
Bahkan, performansi \textit{supertagger} yang dikembangkan sangat baik.
Riset tersebut pada intinya membuktikan bahwasannya MaxEnt dapat digunakan di \textit{supertagger}
juga meskipun pada peruntukannya MaxEnt dibuat untuk POS \textit{tagger}.
Tentunya terdapat beberapa penyesuaian yang harus dilakukan.
Salah satunya adalah formula probabiltas yang digunakan.
Formula yang telah disesuaikan tersebut dapat dilihat di persamaan \ref{maxent:equation:2}.

%

%%%%%%%%%%%%%%%%%%%%%%%%%%%%%%%%%%%%%%%%%%%%%%%%%%%%%%%%%%%%%%%
% Berikut adalah untuk membuat Daftar pustaka
% Style daftar pustaka dapat disesuaikan dengan mengubah 
%\bibliographystyle{kode}
% dengan kode = acm maka hasinya contoh [1]
% dengan kode = agsm maka hasilnya Harvard style (Sukimin, 2017)
%%%%%%%%%%%%%%%%%%%%%%%%%%%%%%%%%%%%%%%%%%%%%%%%%%%%%%%%%%%%%%%%
\cleardoublepage
\addcontentsline{toc}{chapter}{Daftar Pustaka}
\bibliographystyle{plain} %harvard style
\bibliography{References}
%
%\pagebreak
%%%%%%%%%%%%%%%%%%%%%%%%%%%%%%%%%%%%%%%%%%%%%%%%%%%%%%%%%%%%%%%
% Berikut untuk lampiran
%%%%%%%%%%%%%%%%%%%%%%%%%%%%%%%%%%%%%%%%%%%%%%%%%%%%%%%%%%%%%%
\cleardoublepage
\addcontentsline{toc}{chapter}{Lampiran}
\chapter*{Lampiran}

\end{document}
